% Part: lambda-calculus
% Chapter: introduction
% Section: composition

\documentclass[../../../include/open-logic-section]{subfiles}

\begin{document}

\olfileid{lam}{int}{com}
\olsection{\usetoken{S}{lambda definable} फलने संयोजनाच्या अंतर्गत बंद असतात}

\begin{lem}
!!{lambda definable} फलनांचा वर्ग संयोजनाच्या अंतर्गत बंद असतो.
\end{lem}

\begin{proof}
$f$ हे $h$, $g_0,$ \dots,~$g_{k-1}$ यांच्या संयोजनाने
परिभाषित आहे असे समजा. $h$, $g_0$, \dots,~$g_{k-1}$
ही फलने अनुक्रमे $H$, $G_0$, \dots,~$G_{k-1}$
यांनी !!{lambda defined} आहेत असे गृहीत धरल्यास,
~$f$ ला !!{lambda define}s असे पद~$F$ शोधायचे आहे.
पण~$F$ ची व्याख्या सरळ अशी करता येते:
\[
F(x_0, \dots, x_{l-1}) = H(G_0(x_0, \dots, x_{l-1}),
\dots, G_{k-1}(x_0, \dots, x_{l-1})).
\]
दुसऱ्या शब्दांत, लॅम्डा कलनाची भाषा संयोजनाचे निरूपण
करण्यासाठी योग्य आहे.
\end{proof}

\end{document}
